%% kapa: exam=1 round=30 year=2019 period=1 form=A subject=real_estate subject_ko=부동산학원론 seq=20 number=100
%% source: Q-Net 감정평가사 기출문제 (cst00309) artlSeq=5205205 file=extracted/1008/5205205/2019년 제30회 감정평가사 1차시험 문제지/2019년 제30회 감정평가사 1차시험 문제지/2019년도 제30회 감정평가사 1차 1교시 A형.hwp
%% conversion: kordoc 4.24.51 -> tools/gen1.py (auto); review=proofread-sonnet; flags=none
\begin{question}{100}
감정평가에 관한 규칙상 가치에 관한 설명으로 옳지 않은 것은?
\begin{choices}
  \item 대상물건에 대한 감정평가액은 시장가치를 기준으로 결정하는 것을 원칙으로 한다.
  \item 법령에 다른 규정이 있는 경우에는 시장가치 외의 가치를 기준으로 감정평가 할 수 있다.
  \item 대상물건의 특성에 비추어 사회통념상 필요하다고 인정되는 경우에는 시장가치 외의 가치를 기준으로 감정평가 할 수 있다.
  \item 시장가치란 대상 물건이 통상적인 시장에서 충분한 기간 방매된 후 매수인에 의해 제시된 것 중에서 가장 높은 가격을 말한다.
  \item 감정평가 의뢰인이 요청하여 시장가치 외의 가치로 감정평가하는 경우에는 해당 시장가치 외의 가치의 성격과 특징을 검토하여야 한다.
\end{choices}
\end{question}
\answer{4}
